\section{Corrections to the pinned manuscript}\label{sec:corrections}
% Remaining localized corrections found in current consolidated manuscript.
\subsection{The mixed Gaussian antisymmetry must also swap the colors}

The paragraph following the remark about the weight-six diagonal in
\texttt{chapters/04-depth.tex} says that all depth-two content lives in
$\Li_{a,b}(P)-\Li_{b,a}(P)$, ``which vanish on the diagonal.''
This is correct for a point $P=(x,x)$, but not for the mixed points.
Stuffle exchanges both the indices and the colors. The correct pair is
\[
 A_{a,b}(x,y)=\Li_{a,b}(x,y)-\Li_{b,a}(y,x).
\]
Its symmetric partner is a product of singles minus $\Li_{a+b}(xy)$.
For $a=b$, $A_{a,a}(x,y)$ need not vanish when $x\ne y$.

Here is an exact counterexample entirely inside the manuscript's
alphabet:
\begin{equation}\label{audit:eq:mixed-counterexample}
 \Im\bigl(\Li_{1,1}(-1,i)-\Li_{1,1}(i,-1)\bigr)
 =\frac\pi2\log2-G
 =\int_0^1\frac{\log(1+t^2)}{1+t^2}\,dt>0.
\end{equation}
To prove the first equality, differentiating the convergent defining
series in the disk gives
\[
 \frac{\partial}{\partial x}\Li_{1,1}(x,y)
   =\frac{\Li_1(xy)}{1-x}.
\]
Continuation along the segment to $x=-1$ is nonsingular, and hence
\[
 \Im\Li_{1,1}(-1,i)=\int_0^1\frac{\arctan t}{1+t}\,dt.
\]
Put $K=\int_0^1\log(1+t)/(1+t^2)\,dt$. With $t=\tan\theta$,
\[
 K=\int_0^{\pi/4}\left(
   \tfrac12\log2+\log\sin(\theta+\pi/4)-\log\cos\theta
   \right)d\theta=\frac\pi8\log2.
\]
The last two integrals cancel after reflection about $\pi/2$.
Integration by parts now gives
$\Im\Li_{1,1}(-1,i)=\pi\log2/4-K=\pi\log2/8$.
Stuffle and
\[
 \Li_1(-1)=-\log2,\qquad
 \Li_1(i)=-\tfrac12\log2+\tfrac{\pi i}{4},\qquad
 \Im\Li_2(-i)=-G
\]
give
\[
 \Im\bigl(\Li_{1,1}(-1,i)+\Li_{1,1}(i,-1)\bigr)
   =G-\frac\pi4\log2.
\]
Subtracting proves the first equality in
\eqref{audit:eq:mixed-counterexample}; the second also follows directly by setting $t=\tan\theta$. Indeed,
\[
 \int_0^{\pi/4}\log\cos\theta\,d\theta
 =-\frac\pi4\log2+\frac G2,
\]
which follows from the sum and difference of the corresponding
log-sine and log-cosine integrals: their difference is $G$, and their
sum is $-\pi\log2/2$ by the double-angle formula. Thus the displayed
integral is $-2\int_0^{\pi/4}\log\cos\theta\,d\theta$.
Its integrand is positive away from zero. The difference is approximately
$0.17282745097458205019574093418642286289514247590297$.

\paragraph{Suggested replacement.}
``After the symmetric stuffle combinations have been removed, the
remaining information is carried by
$\Li_{a,b}(x,y)-\Li_{b,a}(y,x)$. These vanish when both the indices and
the colors coincide. At the mixed Gaussian points, equality of the
indices alone does not eliminate the antisymmetric combination.''

\subsection{An exact check on the weight-six diagonal example}

The preceding remark in \texttt{04-depth.tex} says that a kernel returns
the diagonal $\Li_{3,3}(i,i)$ in $\beta(4)$, $\pi^4$, and $\zeta(3)$.
This description obscures the simple homogeneous exact formula supplied
by the stated stuffle identity:
\[
 \boxed{\displaystyle
 \Li_{3,3}(i,i)=\frac{9}{2048}\zeta(3)^2
       +\frac{47}{1935360}\pi^6
       -\frac{3i}{1024}\pi^3\zeta(3).}
\]
Indeed, $\Li_3(i)=(-3\zeta(3)+i\pi^3)/32$ and
$\Li_6(-1)=-31\pi^6/30240$, so this is exactly
$\tfrac12(\Li_3(i)^2-\Li_6(-1))$.
No assertion about a former symbolic kernel's output is needed.

The nearby sentence saying that every real single value is a rational
multiple of $\zeta(n)$ also needs the $n=1$ exception:
$\Re\Li_1(i)=-\log2/2$ and $\Li_1(-1)=-\log2$.
For $n\geq2$ the displayed zeta description is correct.
The analogous Eisenstein sentence also requires $n\ge2$; its separate
order-one value is $\Li_1(e^{2\pi i/3})=-\log3/2+i\pi/6$.

\subsection{The known component of a digamma difference}

In \texttt{chapters/06-cm.tex}, lines 91--93 call the real component at
odd weight and imaginary component at even weight ``opaque.'' The
preceding two displayed formulas evaluate exactly those components.
This is a parity reversal in the commentary, rather than an error in
those displayed identities.

For $0<c<1$, put $z=(1-c)/2+it$ and
\[
 R_c(t)=\frac{\pi\sin(\pi c)}{\cosh(2\pi t)+\cos(\pi c)}.
\]
Reflection and conjugation give
$\Re(\psi(z+c)-\psi(z))=R_c(t)$. Differentiating $m$ times in the real
variable $t$ proves the exact all-order statement
\[
 \Re\left(i^m\bigl(\psi^{(m)}(z+c)-\psi^{(m)}(z)\bigr)\right)
       =R_c^{(m)}(t).
\]
Thus the real component is determined when $m$ is even, and the imaginary
component when $m$ is odd. Under the manuscript's weight convention
$w=m+1$, these are the real component at odd weight and imaginary component
at even weight. The complementary components are not determined by this
identity. No assertion of nonreducibility follows.

The corrected sentence should read: ``The complementary component---the
imaginary part at odd weight and the real part at even weight---is not
determined by this reflection identity.'' This statement concerns shifted
differences and should not be conflated with the parity of individual
polygamma values on the line $\Re z=1/2$.
